A statistically convincing result can be produced by the fitting procedure rather than the
data, and no further statistical testing will detect it. We document four instances in one
pre-registered study of whether anchoring a quantile-loss model to a Value-at-Risk
prior improves one-day-ahead forecasts. It does not; that answer is not the contribution.
\textbf{All four findings were withdrawn} --- one to a multiplicity correction, one to a power
analysis, one to the model's own seed noise. The fourth, replicated across four runs with a
dose--response and a sign test at \(p = 5.2 \times 10^{-4}\), dissolved when asked whether the
mechanism guaranteed it: 28 minutes of compute, no test data.

\textbf{Neither the mechanism (ridge shrinkage) nor the control (a scale-matched permutation) is
new}; the artefact is reportable and survives pre-registration. \textbf{The market
study is the exhibit, not the evidence:} we disclose \textbf{3,555 specifications evaluated} and
\textbf{1,959 test-set evaluations}, zero cells scored once; every market result is
validation-grade. The evidence is an analytic derivation and a synthetic reproduction with
ground truth known: an uninformative anchor stabilises the estimator as much as the true
optimum (\textbf{median ratio 1.04, 95\% CI 1.0--1.8}) while forecasting worse. \textbf{Stability tracks the
penalty; usefulness tracks the target.}

\begin{center}\rule{0.5\linewidth}{0.5pt}\end{center}
